% Reconstructed from the four exact author transmissions retained by root.
% Root and B: full reconstruction/source audit PASS.
% D: independent targeted Artinian pole and Isom-openness audit PASS.
\begin{lemma}[Proved: Polynomial witnesses for bounded derived diagrams]
\label{D:polynomial-monad-diagrams}
We fix a projective space $P=\mathbf P^N_k$, finitely many vector bundles
on $P$, a bound on complex length, and a finite list of operations
consisting of tensoring with fixed bounded complexes, absolute Frobenius
pullback, composition, and equality of derived morphisms. Suppose the
variable complexes have at most $cr$ copies of each fixed bundle in each
degree. A fixed finite diagram formed from these operations admits a
witness scheme defined by at most $Cr^2$ scalar polynomial equations in
at most $Cr^2$ variables, of degrees at most $D$. The constants depend
on the fixed data and on $p$, but not on $r$. The equations interpret
the diagram correctly over every $k$-algebra; they parameterize
representatives and witnesses for equality, rather than a quotient of
the space of representatives.
\end{lemma}
\begin{proof}
We choose the standard finite affine cover of $P$. For each pair of
fixed bundles, its Cech Hom complex is a complex of $k$-vector spaces
with finite-dimensional cohomology. We split it $k$-linearly into its
cohomology and contractible summands. For two variable bounded complexes,
the total Hom differential has the form $d_0+\delta$, where $d_0$ is
the fixed Cech differential and $\delta$ raises the complex-Hom column.
The column has a fixed finite range. Consequently the elimination of
the contractible summands uses a finite geometric series in $h\delta$,
where $h$ is the chosen contraction. Every resulting differential,
inclusion, projection and homotopy is a polynomial of bounded degree
in the differential coordinates. The finite remaining complex has
dimension at most $Cr^2$: its terms are sums of fixed cohomology spaces
tensored with matrix spaces of dimensions at most $(cr)^2$.

We impose the cycle equation on a degree-zero representative. An equality
of derived morphisms is imposed by adjoining a degree-minus-one variable
whose differential is the difference. Composition is computed by the
Cech cup product and then the finite projection. Tensoring with a fixed
bounded complex enlarges the fixed bundle list and the column range by
fixed amounts. Frobenius replaces coefficients by their $p$th powers
and replaces each fixed bundle by its fixed Frobenius pullback. The
finite list of operations therefore preserves the asserted bounds.
All splittings are over $k$, so their identities remain identities
after tensoring with any $k$-algebra. In particular this calculation
applies to nonreduced Artinian parameter rings. We do not replace the
derived category by its Ext algebra or assume that this algebra has a
strict finite-dimensional differential graded model.
\end{proof}

\begin{theorem}[Proved: A quadratic rank bound for defining-characteristic
representations on a punctured proper surface]
\label{D:global-quadratic-representation-bound}
Let $S$ be a normal proper connected surface over an algebraically closed
field $k$ of characteristic $p>0$, and let $T\subset S$ be a finite closed
subset. We put $U=S\setminus T$ and assume $U\ne\varnothing$. There is
a constant $C_S$ such that, for every $r,f\ge1$, with $q=p^f$,
\[
 \#\bigl(\operatorname{Hom}_{\mathrm{cont}}
   (\pi_1(U),\operatorname{GL}_r(\mathbf F_q))/
   \operatorname{GL}_r(\mathbf F_q)\bigr)
 \le q^{C_Sr^2}.
\]
The same conclusion, with a constant depending on the scheme and boundary,
holds for a proper scheme of dimension at most two whose finite removed
subset has codimension at least two in every irreducible component.
This statement concerns the defining coefficient characteristic. It does
not assert finite generation of the full fundamental group.
\end{theorem}
\begin{proof}
We first work on $S_{\mathrm{reg}}$. We fix a smooth projective resolution
$\pi:Y\to S$, an effective exceptional cycle $Z$ with $Z\cdot E<0$ for
every exceptional component $E$, and an embedding
$i:Y\hookrightarrow P=\mathbf P^N_k$. If the exceptional locus is empty,
we set $Z=0$. We write $V=Y\setminus\operatorname{Supp}Z$.

\emph{Uniform linear sizes.}
A rank-$r$ $\mathbf F_{p^f}$ local system gives reflexive
coefficient-Frobenius stages $E_a$, indexed by
$a\in\mathbf Z/f\mathbf Z$, and isomorphisms
$F_S^{[*]}E_a\simeq E_{a+1}$. We choose their minimal extensions
$F_a$ to $Y$. The proof of
Theorem~\ref{D:global-periodic-reflexive-boundedness} gives the following
bounds uniformly in $r$, $f$ and $a$: fixed upper and lower slope windows
on each exceptional component; $c_1(F_a)$ in the exceptional span with
coefficients $O(r)$; $\chi(F_a)=O(r)$; and fixed upper bounds for the
slopes of $F_a$ and its dual on $Y$ and on one fixed smooth ample curve.
In particular, for every fixed integer $t$,
\[
 h^j(Y,F_a(t))\le C_tr\qquad(0\le j\le2).
\]
Here the bounds for $h^0$ and $h^2$ follow from the fixed-curve restriction
and duality arguments already proved, and Riemann--Roch supplies the bound
for $h^1$. The same conclusion holds after tensoring by any one of
finitely many fixed vector bundles: we take a fixed three-step partial
resolution by sums of fixed line twists, tensor by $F_a$, and use the
vanishing of coherent cohomology in degrees above two. It also holds for
$F_Y^*F_a$ and its finitely many fixed twists. Indeed the one-step pole
comparison bounds its slopes and dual slopes; Frobenius scales the first
and second Chern characters by $p$ and $p^2$, respectively. We use only
this fixed one-step Frobenius, not a complex whose degree grows with $p^f$.

By the uniform pole lemma we may choose one integer $B$ such that the
isomorphisms on $V$ between $F_Y^*F_a$ and $F_{a+1}$, and their inverses,
extend with poles at most $BZ$. We enlarge $B$ once to bound maps in both
directions between any two bundles in the fixed exceptional slope windows.
The same bound will later apply to isomorphisms over Artinian parameter
rings. Put $L=\mathcal O_Y(BZ)$ and let $s$ be its canonical section.

The Beilinson monad of a coherent sheaf $\mathcal F$ on $P$ has terms
\[
 \mathcal B^e(\mathcal F)=
 \bigoplus_j H^j(P,\mathcal F(e-j))\otimes\Omega_P^{j-e}(j-e),
\]
where $0\le j-e\le N$; it is exact except in degree zero, where its
cohomology is $\mathcal F$
\citep[Theorem~6.1, pp.~23--25]{C_EFS2003}. Applied to the pushforward
of any of the bundles just listed, only $j=0,1,2$ occurs. Thus there is a
fixed finite set of bundle types and degrees, and every multiplicity is
at most $Kr$, for one fixed $K$. Per Frobenius stage the number of possible
multiplicity arrays is at most $(Kr+1)^{K_0}$, with $K_0$ fixed.

\emph{The parameter schemes.}
We describe them explicitly in order to keep their degree and size bounds
independent of the rank. Choose fixed finite vector-bundle resolutions
$K_L$ and $K_0$ of $i_*L$ and $i_*\mathcal O_Y$ on $P$. These exist by
the finite graded free resolution over the polynomial coordinate ring.
Fix also a derived representative of the section
$s:\mathcal O_P\to i_*L$.

For each $a$ we use six monads, representing
\[
 F_a,\quad T(F_a),\quad T^2(F_a),\quad
 P_a,\quad T(P_a),\quad T^2(P_a),\qquad
 T(M)=M\otimes L,\quad P_a=F_Y^*F_a.
\]
The $T$-nodes are introduced iteratively. Given a monad $C_M$, we introduce
a monad $C_{T(M)}$ and a derived map
\[
 \beta_M:C_M\otimes K_L\longrightarrow C_{T(M)}.
\]
We require its cone to have fibre cohomology only in degrees at most $-2$.
Once the monads have only degree-zero cohomology, this condition identifies
$H^0(C_{T(M)})$ with $H^0(C_M)\otimes L$. For the $P_a$-node we use in
the same way
\[
 LF_P^*C_{F_a}\otimes K_0\longrightarrow C_{P_a}.
\]
The naive pullback of the vector-bundle complex computes this derived
pullback. Absolute Frobenius need not be flat over a nonreduced parameter
ring; all its negative Tor terms, as well as those from tensoring with
$K_0$ or $K_L$, are retained. Their degree-zero cohomology is the desired
ordinary pullback or tensor product. The cone condition is precisely what
identifies that cohomology with the auxiliary monad.

The fixed section $s$ gives maps $\sigma_M:M\to T(M)$. We introduce four
arrows per stage,
\[
 a_a:P_a\to T(F_{a+1}),\qquad
 b_a:F_{a+1}\to T(P_a),\qquad T(a_a),\qquad T(b_a).
\]
The last two arrows are forced to be the tensor transforms of the first
two by the equations
\[
 T(a_a)\beta_{P_a}
   =\beta_{T(F_{a+1})}(a_a\otimes K_L),\qquad
 T(b_a)\beta_{F_{a+1}}
   =\beta_{T(P_a)}(b_a\otimes K_L).
\]
We impose
\[
 T(b_a)a_a=\sigma_{T(P_a)}\sigma_{P_a},\qquad
 T(a_a)b_a=\sigma_{T(F_{a+1})}\sigma_{F_{a+1}}.
\]
On $V$, division by $s$ makes $a_a$ and $b_a$ inverse isomorphisms.
This iterative six-node diagram makes the tensor identifications
coherent; we do not choose independent, unconnected identifications
for the second twists.

All equations just described are imposed using
Lemma~\ref{D:polynomial-monad-diagrams}. The fixed bundle list consists of
the $\Omega_P^j(j)$, their finitely many twists by the degrees of generators
of $I_Y$, the terms $\Omega_P^j(j)\otimes K_L^b$, and the terms
$F_P^*\Omega_P^j(j)\otimes K_0^b$. The complex lengths are fixed. At each
stage there are six monads, five auxiliary degree-zero identification
maps, four arrow variables, and a fixed number of derived equalities.
Each differential block, cycle space and boundary-witness space has
dimension $O(r^2)$. We consequently obtain at most $Afr^2$ variables
and at most $Afr^2$ scalar equations, all of degrees at most $D$, for
constants $A,D$ depending only on the fixed geometry and $p$.

We impose the scheme-theoretic support condition by choosing fixed
homogeneous generators $g$ of $I_Y$ and requiring multiplication
$g:C(-\deg g)\to C$ to be zero as a derived morphism. A boundary witness
expresses this as polynomial equations. We then take the open locus on
which every monad has cohomology only in degree zero and that cohomology
is a base-flat sheaf $i_*F$, with $F$ locally free of rank $r$ on $Y$.
The relative truncation lemma of Lieblich gives this openness and
compatibility with base change
\citep[Lemma~2.1.4, pp.~4--5]{C_Lieblich2006}. Local freeness follows
from flatness over the parameter base and local freeness on its fibres:
we lift a fibre basis, and apply Nakayama to the cokernel and then, using
flatness, to the kernel. Properness of $Y$ makes this an open condition
on the parameter base. We also impose the open fixed slope-window
conditions on the original $F_a$, and the open cone conditions above.
The slope conditions may be expressed using the proper images of the
finitely many relevant Quot schemes of coherent quotients; torsion can
only increase their degrees, so including all coherent quotients gives
the correct closed bad loci.

We denote the resulting open witness scheme for the chosen arrays by $W$.
It has a universal family of bundles on $Y$ and a universal unit Frobenius
cycle on $V$. Every local system under consideration has a point in one
of these schemes: its six bundles have the stated linear monad bounds,
and the actual tensor maps, inverse arrows and derived equalities provide
all the witnesses. We need no quantitative bound for the complexity of
the open conditions. An open subset of an irreducible scheme remains
irreducible when nonempty.

\emph{Uniform poles over Artinian bases.}
We spell out this point because the bundles chosen on $Y$ can vary even
when their restrictions to $V$ are isomorphic. If $F,G$ lie in the fixed
exceptional slope windows, a fixed filtration of the nonreduced cycle
$Z$ by line-bundle twists on its reduced components gives
\[
 H^0\bigl(Z,\mathcal Hom(F,G)\otimes\mathcal O_Z(nZ)\bigr)=0
 \qquad(n>B),
\]
with $B$ independent of the ranks. Indeed $Z\cdot E<0$, and on each fixed
layer the maximum target slope eventually lies below the minimum source
slope. This uses the Hom vanishing criterion, not semistability of a
tensor product.

Now let $A$ be any local Artinian $k$-algebra and let $F_A,G_A$ be two
families whose closed fibres lie in those windows. The Hom bundle
restricted to each fixed $Z$-layer is flat over $A$. Its maximal-ideal
filtration has graded pieces equal to the closed-fibre sheaf tensored
with vector spaces. The same $H^0$ vanishing therefore holds over $A$.
From
\[
 0\to\mathcal Hom(F_A,G_A)((n-1)Z)
 \to\mathcal Hom(F_A,G_A)(nZ)
 \to\mathcal Hom(F_A,G_A)\otimes\mathcal O_Z(nZ)\to0
\]
we see that the global Hom spaces stabilize at $B$. Every map on $V_A$
first extends with some finite pole order, because localization along
the effective Cartier divisor $Z$ is the filtered colimit of these
twists and global sections commute with this colimit on the
quasi-compact separated scheme $Y_A$. It therefore extends with poles
at most $BZ$. The same holds for its inverse. This bound is independent
of the length of $A$, of $r$, and of $f$. Multiplication by a local
equation of $Z$ is injective on these bundles, even over nonreduced $A$,
so the extension is unique.

\emph{Open unit-isomorphism classes.}
Fix a closed point $w\in W$ and its reference unit cycle. Consider the
functor of isomorphisms from this reference cycle to the variable cycle
on $V$, represented by maps in both directions on $Y$ with poles at most
$BZ$. The proper Hom functor is affine of finite presentation
\citep[Proposition~99.3.10, Tag~08K6]{C_StacksQuantitativeMonads}.
The inverse equations and the one-step Frobenius compatibility equations,
after clearing fixed poles, define an affine finitely presented unit-Isom
scheme over $W$. The compatibility equations involve $pB$ and $B$,
rather than $p^fB$, because we retain all $f$ consecutive $p$-stages.

For a small extension of local Artinian parameter algebras
$A\twoheadrightarrow A/I$, a specified unit-isomorphism over $A/I$
lifts uniquely on $V_A$. Indeed a unit $p^f$-Frobenius vector bundle has
a finite etale $\operatorname{GL}_r(\mathbf F_{p^f})$ frame torsor,
and finite etale covers and their morphisms are invariant under
nilpotent thickenings
\citep[Lemma~58.8.3, Tag~0BQB]{C_StacksQuantitativeMonads}. The frame
torsor can also be checked directly from the equations for
Frobenius-fixed frames: their derivative is invertible, and over every
geometric point they give the usual $\mathbf F_{p^f}$ basis torsor.
The preceding uniform pole argument extends the lifted isomorphism and
its inverse to $Y_A$ with the same bound $B$. Their equations hold
globally because they hold on $V_A$ and localization along $Z$ is
injective. Thus the unit-Isom scheme lifts the specified Artinian point.
The infinitesimal smoothness criterion applies
\citep[Lemma~10.141.2, Tag~02HT]{C_StacksQuantitativeMonads}; its image
contains an open neighbourhood of $w$ consisting of its unit-isomorphism
class.

We repeat this at every closed point of $W$. Distinct unit-isomorphism
classes give disjoint open subsets: an intersection, if nonempty, would
contain a closed point because $W$ is of finite type over $k$. These opens
cover $W$, since their complement is closed and has no closed point.
Quasi-compactness then gives a finite clopen partition. Therefore the
number of classes represented by $W$ is at most its number of connected
components, and hence at most its number of irreducible components.
This argument does not infer algebraicity from formal constancy of
covers on an arbitrary nonproper scheme. The fixed exceptional pole
bound is what makes the unit-Isom functor finitely presented.

An affine scheme cut out by $m$ polynomials of degree at most $D\ge2$
has at most $D^m$ irreducible components. For completeness, we homogenize
the equations and intersect them successively in projective space.
An irreducible component contained in the next hypersurface retains its
degree; for any other component the sum of degrees of its
hypersurface-section components is at most $D$ times its degree, by
the multiplication exact sequence for the hypersurface equation and
its Hilbert polynomial. Induction bounds the sum of degrees, and thus
the number of irreducible components, by $D^m$. Passing to the affine
chart and then to an open subset cannot increase that number. Here
$m\le Afr^2$. The multiplicity-array choices contribute at most
$(Kr+1)^{K_0f}$. Absorbing this factor gives
\[
 \#\{\text{represented unit classes}\}\le\exp(Cfr^2)
 =p^{C'fr^2}.
\]
Unit cycles up to unit-isomorphism are exactly the $\mathbf F_{p^f}$
local systems. This proves the bound on $S_{\mathrm{reg}}$. Removing
the regular points of $T$ does not change its fundamental group by
purity. Restriction from $U$ to $S_{\mathrm{reg}}\setminus T$ is fully
faithful on connected finite covers, since $U$ is normal and the latter
open is dense. Hence its fundamental group surjects onto $\pi_1(U)$,
and the same bound applies to $U$.

Finally, for a possibly nonnormal proper scheme of dimension at most two,
reduction does not change finite etale covers. Its normalization is
finite. On each normalized surface component we use the bound just
proved; proper curve and point components have the analogous bound
from their finitely generated fundamental groups. For fixed pulled-back
local systems, effective finite-surjective descent identifies descent
data with isomorphisms on the double fibre product satisfying the
cocycle equation
\citep[Expose~IX, Theorem~4.7]{C_SGA1}. A section of the resulting Isom
local system is determined by its value at one geometric point of each
connected component. Thus the number of possible descent data is at
most $|\operatorname{GL}_r(\mathbf F_q)|^c\le q^{cr^2}$, for the fixed
finite number $c$ of double-product components. Multiplying these bounds
proves the stated extension. We treat each connected component separately
when necessary.
\end{proof}
